\documentclass[../main.tex]{subfiles}
\begin{document}

\begin{center}
    \Large{\textbf{ABSTRACT}}\\
\end{center}
\vspace{1cm}
In the context of the rapidly growing e-commerce sector in Vietnam, small and medium-sized enterprises (SMEs) frequently encounter significant financial and technical barriers when attempting to build and operate independent online stores. Existing e-commerce platforms on the market, such as Shopify, WooCommerce, or Haravan, still exhibit considerable limitations. Specifically, they fail to simultaneously satisfy three crucial criteria: low initial and maintenance costs, flexible interface customization without requiring programming knowledge, and system scalability capable of serving tens of thousands of stores without linearly increasing infrastructure costs. Stemming from this practical need, this thesis focuses on the research and development of a multi-tenant e-commerce platform named ShopVolo v2 to thoroughly address the resource optimization problem. The primary approach of the study involves implementing a Multi-tenant Software as a Service (SaaS) architecture combined with the "Zero-file" design principle, wherein the entire storefront interface is represented as configuration data rather than discrete physical source code files.

The system is developed based on a Monorepo model using the Turborepo tool, integrating modern technology stacks including Next.js for the administrative dashboard interface and dynamic storefront rendering, and NestJS for the core API system with automatic user data isolation capabilities. Furthermore, a dual database system combining PostgreSQL and MongoDB is employed to efficiently process both structured information and flexible configuration documents. The highlight of the system architecture is the Deep Merge algorithm, which allows the real-time blending of the standard interface structure with individual custom configurations, combined with Redis caching to ensure near-instantaneous response speeds. [TODO: Supplement details regarding the main contributions and specific performance results achieved after the deployment and testing of the system].

\end{document}